\documentclass[11pt,a4paper]{article}
\usepackage[utf8]{inputenc}
\usepackage[T1]{fontenc}
\usepackage[catalan]{babel}
\usepackage{lmodern}
\usepackage[sfdefault,lining]{FiraSans}
\usepackage[a4paper,top=2.3cm,bottom=2.2cm,left=1.8cm,right=1.8cm,headheight=1.3cm,headsep=0.4cm,footskip=1cm]{geometry}
\usepackage{amsmath,amssymb}
\usepackage{icomma}
\usepackage{graphicx}
\usepackage[table]{xcolor}
\usepackage{tikz}
\usepackage{fontawesome5}
\usepackage{ccicons}
\usepackage{fancyhdr}
\usepackage{lastpage}
\usepackage{tasks}
\usepackage{needspace}
\usepackage[most]{tcolorbox}
\usepackage{array}
\usepackage[hidelinks,pdftitle={Solucions de la fitxa de la unitat 1: Nombres naturals, enters i fraccionaris. 4t d'ESO Matemàtiques A},pdfsubject={Matemàtiques A 4t ESO},pdfkeywords={nombres naturals, operacions combinades, potències, arrels, nombres enters, valor absolut, fraccions, fraccions equivalents, simplificació, amplificació, problemes amb fraccions, propietats de les potències}]{hyperref}

\definecolor{navy}{HTML}{1B2748}
\definecolor{accent}{HTML}{E8603A}
\definecolor{teal}{HTML}{1F9E90}
\definecolor{soft}{HTML}{EEF1F8}
\definecolor{muted}{HTML}{5D6685}

\setlength{\parindent}{0pt}
\setlength{\parskip}{3pt}
\settasks{label=\alph*), label-format=\bfseries\color{navy}, label-width=1.3em, item-indent=1.8em, column-sep=0.6em, after-item-skip=2pt, before-skip=3pt, after-skip=4pt}

% ---------- capçalera i peu ----------
\newcommand{\llicencia}{\href{https://creativecommons.org/licenses/by-nc-sa/4.0/deed.ca}{\ccbyncsa}\ \ Llicència Creative Commons Reconeixement-NoComercial-CompartirIgual 4.0 Internacional (CC BY-NC-SA 4.0)}
\pagestyle{fancy}
\fancyhf{}
\fancyhead[L]{\raisebox{-0.35\height}{\includegraphics[height=0.95cm]{logo-gva.png}}}
\fancyhead[C]{\small\color{muted}Solucions de la fitxa · Unitat 1: Nombres naturals, enters i fraccionaris · 4t ESO A}
\fancyhead[R]{\raisebox{-0.35\height}{\includegraphics[height=0.95cm]{logo-vermellar.png}}}
\fancyfoot[L]{\scriptsize\color{muted}\llicencia}
\fancyfoot[R]{\small\color{muted}\thepage}
\renewcommand{\headrulewidth}{0.4pt}
\renewcommand{\headrule}{\color{muted!40}\hrule height \headrulewidth}
\fancypagestyle{primera}{\fancyhf{}\renewcommand{\headrulewidth}{0pt}
  \fancyfoot[L]{\scriptsize\color{muted}\llicencia}
  \fancyfoot[R]{\small\color{muted}\thepage}}

% ---------- nivells i exercicis ----------
\newcommand{\estrelles}[1]{{\color{accent}\ifcase#1\or\faStar\faStar[regular]\faStar[regular]\or\faStar\faStar\faStar[regular]\or\faStar\faStar\faStar\fi}}
\newcounter{exer}
\newcommand{\ex}[1]{\par\needspace{4\baselineskip}\medskip\stepcounter{exer}\noindent{\color{navy}\bfseries\theexer.}\ \ {\small\estrelles{#1}}\ \ }

\newcommand{\seccio}[2]{\par\needspace{10\baselineskip}\bigskip
  \begin{tcolorbox}[colback=navy,colframe=navy,coltext=white,arc=2mm,boxrule=0pt,left=3mm,right=3mm,top=1.2mm,bottom=1.2mm]
  {\large\bfseries #1\quad #2}
  \end{tcolorbox}\vspace{-2pt}}

\newcolumntype{C}[1]{>{\centering\arraybackslash}p{#1}}

% quadre de recordatori (només a la fitxa d'exercicis)
\newcommand{\recorda}[1]{\par\smallskip\begin{tcolorbox}[colback=accent!7,colframe=accent!55,arc=2mm,boxrule=0.5pt,left=2.5mm,right=2.5mm,top=1mm,bottom=1mm]\small{\color{accent}\faLightbulb[regular]}\ \textbf{Recorda:}\ #1\end{tcolorbox}}
\newcommand{\res}[1]{\ifmmode\boxed{#1}\else\fbox{#1}\fi}

\def\fmneed{11}% que el títol d'una secció no quede sol al final del full
% =====================================================================
% Capçalera FentMates de les fitxes: un color per a cada curs.
% Els colors són els de la pàgina inicial (LEVELS a assets/js/materials.js);
% el logo ve de _eines/logo/pdf/ (python3 _eines/logo/genera.py).
%
% Ús, al final del preàmbul de la unitat:
%   % =====================================================================
% Capçalera FentMates de les fitxes: un color per a cada curs.
% Els colors són els de la pàgina inicial (LEVELS a assets/js/materials.js);
% el logo ve de _eines/logo/pdf/ (python3 _eines/logo/genera.py).
%
% Ús, al final del preàmbul de la unitat:
%   % =====================================================================
% Capçalera FentMates de les fitxes: un color per a cada curs.
% Els colors són els de la pàgina inicial (LEVELS a assets/js/materials.js);
% el logo ve de _eines/logo/pdf/ (python3 _eines/logo/genera.py).
%
% Ús, al final del preàmbul de la unitat:
%   \input{../comu/fentmates.tex}
%   \fmcurs{1eso}
% i, al començament del document:
%   \capcalera{Matemàtiques}{1r ESO}{Fitxa d'exercicis · Unitat 2}{Divisibilitat}
% Opcional abans de \input: \def\fmneed{14} (espai mínim abans de cada secció).
% =====================================================================
\makeatletter
\def\fm@logodir{../../logo/pdf}
\providecommand{\fmneed}{10}
\def\fm@cursnom{}

% clau del curs, color i color del text damunt del color
\newcommand{\fm@curs}[3]{\@namedef{fm@color@#1}{#2}\@namedef{fm@text@#1}{#3}}
\fm@curs{1eso}{E6B84A}{navy}
\fm@curs{2eso}{3B82F6}{white}
\fm@curs{3eso}{FF7A50}{white}
\fm@curs{4eso-a}{A78BFA}{white}
\fm@curs{4eso-b}{6D28D9}{white}
\fm@curs{1batx-mi}{9B1C1C}{white}
\fm@curs{1batx-ccss}{EF6B6B}{white}
\fm@curs{2batx-mii}{1E6B47}{white}
\fm@curs{2batx-ccss}{2BB3A3}{white}

\newcommand{\fmcurs}[1]{%
  \@ifundefined{fm@color@#1}{\PackageError{fentmates}{Curs desconegut: #1}{}}{}%
  \edef\fm@tmp{\noexpand\definecolor{curs}{HTML}{\@nameuse{fm@color@#1}}%
               \noexpand\colorlet{curstext}{\@nameuse{fm@text@#1}}}\fm@tmp
  \colorlet{cursfosc}{curs!72!black}%
  \def\fm@logo{\fm@logodir/vertical-#1.pdf}%
  \def\fm@marca{\fm@logodir/marca-#1.pdf}}

% ---------- capçalera de la primera pàgina ----------
% #1 matèria (petit), #2 curs, #3 tipus de document i unitat, #4 títol
\newcommand{\capcalera}[4]{%
  \gdef\fm@cursnom{#2}%
  \par\noindent
  \begin{tcolorbox}[colback=curs,colframe=curs,boxrule=0pt,arc=2.2mm,width=3.5cm,height=2.15cm,
      valign=center,left=3mm,right=2mm,top=1mm,bottom=1mm,nobeforeafter,box align=center]
    \color{curstext}{\footnotesize #1}\par\vspace{1pt}{\fontsize{19}{21}\selectfont\bfseries #2\par}%
  \end{tcolorbox}\hspace{5mm}%
  \begin{minipage}[c]{\dimexpr\linewidth-3.5cm-5mm-2.9cm\relax}
    {\color{cursfosc}\small\bfseries\MakeUppercase{#3}}\par\vspace{2pt}
    {\color{navy}\fontsize{20}{23}\selectfont\bfseries #4\par}
  \end{minipage}\hfill
  \begin{minipage}[c]{2.6cm}\raggedleft\includegraphics[width=2.6cm]{\fm@logo}\end{minipage}\par
  \vspace{2pt}}

% ---------- barres de secció del color del curs ----------
\renewcommand{\seccio}[2]{\par\needspace{\fmneed\baselineskip}\bigskip
  \begin{tcolorbox}[colback=curs,colframe=curs,coltext=curstext,arc=2mm,boxrule=0pt,left=3mm,right=3mm,top=1.2mm,bottom=1.2mm]
  {\large\bfseries #1\quad #2}
  \end{tcolorbox}\vspace{-2pt}}

% ---------- estrelles de dificultat del color del curs (a Batxillerat no se'n posen) ----------
\@ifundefined{estrelles}{}{%
  \renewcommand{\estrelles}[1]{{\color{curs}\ifcase#1\or\faStar\faStar[regular]\faStar[regular]\or\faStar\faStar\faStar[regular]\or\faStar\faStar\faStar\fi}}}

% ---------- capçalera i peu de les pàgines ----------
% primera pàgina: logos institucionals al peu
\fancypagestyle{primera}{\fancyhf{}\renewcommand{\headrulewidth}{0pt}
  \fancyfoot[L]{\raisebox{-0.32\height}{\includegraphics[height=0.78cm]{logo-gva.png}}\hspace{5mm}%
    \raisebox{-0.32\height}{\includegraphics[height=0.78cm]{logo-vermellar.png}}}
  \fancyfoot[R]{\scriptsize\color{muted}\href{https://creativecommons.org/licenses/by-nc-sa/4.0/deed.ca}{\ccbyncsa}\ \ CC BY-NC-SA 4.0 · fentmates.org\qquad\small\thepage}}
% la resta: el símbol a l'esquerra i el curs a la dreta, amb el títol al centre com abans
\fancyhead[L]{\raisebox{-0.3\height}{\includegraphics[height=0.62cm]{\fm@marca}}}
\fancyhead[R]{\small\bfseries\color{cursfosc}\fm@cursnom}
\renewcommand{\headrule}{\color{curs!45}\hrule height \headrulewidth}
\makeatother

%   \fmcurs{1eso}
% i, al començament del document:
%   \capcalera{Matemàtiques}{1r ESO}{Fitxa d'exercicis · Unitat 2}{Divisibilitat}
% Opcional abans de \input: \def\fmneed{14} (espai mínim abans de cada secció).
% =====================================================================
\makeatletter
\def\fm@logodir{../../logo/pdf}
\providecommand{\fmneed}{10}
\def\fm@cursnom{}

% clau del curs, color i color del text damunt del color
\newcommand{\fm@curs}[3]{\@namedef{fm@color@#1}{#2}\@namedef{fm@text@#1}{#3}}
\fm@curs{1eso}{E6B84A}{navy}
\fm@curs{2eso}{3B82F6}{white}
\fm@curs{3eso}{FF7A50}{white}
\fm@curs{4eso-a}{A78BFA}{white}
\fm@curs{4eso-b}{6D28D9}{white}
\fm@curs{1batx-mi}{9B1C1C}{white}
\fm@curs{1batx-ccss}{EF6B6B}{white}
\fm@curs{2batx-mii}{1E6B47}{white}
\fm@curs{2batx-ccss}{2BB3A3}{white}

\newcommand{\fmcurs}[1]{%
  \@ifundefined{fm@color@#1}{\PackageError{fentmates}{Curs desconegut: #1}{}}{}%
  \edef\fm@tmp{\noexpand\definecolor{curs}{HTML}{\@nameuse{fm@color@#1}}%
               \noexpand\colorlet{curstext}{\@nameuse{fm@text@#1}}}\fm@tmp
  \colorlet{cursfosc}{curs!72!black}%
  \def\fm@logo{\fm@logodir/vertical-#1.pdf}%
  \def\fm@marca{\fm@logodir/marca-#1.pdf}}

% ---------- capçalera de la primera pàgina ----------
% #1 matèria (petit), #2 curs, #3 tipus de document i unitat, #4 títol
\newcommand{\capcalera}[4]{%
  \gdef\fm@cursnom{#2}%
  \par\noindent
  \begin{tcolorbox}[colback=curs,colframe=curs,boxrule=0pt,arc=2.2mm,width=3.5cm,height=2.15cm,
      valign=center,left=3mm,right=2mm,top=1mm,bottom=1mm,nobeforeafter,box align=center]
    \color{curstext}{\footnotesize #1}\par\vspace{1pt}{\fontsize{19}{21}\selectfont\bfseries #2\par}%
  \end{tcolorbox}\hspace{5mm}%
  \begin{minipage}[c]{\dimexpr\linewidth-3.5cm-5mm-2.9cm\relax}
    {\color{cursfosc}\small\bfseries\MakeUppercase{#3}}\par\vspace{2pt}
    {\color{navy}\fontsize{20}{23}\selectfont\bfseries #4\par}
  \end{minipage}\hfill
  \begin{minipage}[c]{2.6cm}\raggedleft\includegraphics[width=2.6cm]{\fm@logo}\end{minipage}\par
  \vspace{2pt}}

% ---------- barres de secció del color del curs ----------
\renewcommand{\seccio}[2]{\par\needspace{\fmneed\baselineskip}\bigskip
  \begin{tcolorbox}[colback=curs,colframe=curs,coltext=curstext,arc=2mm,boxrule=0pt,left=3mm,right=3mm,top=1.2mm,bottom=1.2mm]
  {\large\bfseries #1\quad #2}
  \end{tcolorbox}\vspace{-2pt}}

% ---------- estrelles de dificultat del color del curs (a Batxillerat no se'n posen) ----------
\@ifundefined{estrelles}{}{%
  \renewcommand{\estrelles}[1]{{\color{curs}\ifcase#1\or\faStar\faStar[regular]\faStar[regular]\or\faStar\faStar\faStar[regular]\or\faStar\faStar\faStar\fi}}}

% ---------- capçalera i peu de les pàgines ----------
% primera pàgina: logos institucionals al peu
\fancypagestyle{primera}{\fancyhf{}\renewcommand{\headrulewidth}{0pt}
  \fancyfoot[L]{\raisebox{-0.32\height}{\includegraphics[height=0.78cm]{logo-gva.png}}\hspace{5mm}%
    \raisebox{-0.32\height}{\includegraphics[height=0.78cm]{logo-vermellar.png}}}
  \fancyfoot[R]{\scriptsize\color{muted}\href{https://creativecommons.org/licenses/by-nc-sa/4.0/deed.ca}{\ccbyncsa}\ \ CC BY-NC-SA 4.0 · fentmates.org\qquad\small\thepage}}
% la resta: el símbol a l'esquerra i el curs a la dreta, amb el títol al centre com abans
\fancyhead[L]{\raisebox{-0.3\height}{\includegraphics[height=0.62cm]{\fm@marca}}}
\fancyhead[R]{\small\bfseries\color{cursfosc}\fm@cursnom}
\renewcommand{\headrule}{\color{curs!45}\hrule height \headrulewidth}
\makeatother

%   \fmcurs{1eso}
% i, al començament del document:
%   \capcalera{Matemàtiques}{1r ESO}{Fitxa d'exercicis · Unitat 2}{Divisibilitat}
% Opcional abans de \input: \def\fmneed{14} (espai mínim abans de cada secció).
% =====================================================================
\makeatletter
\def\fm@logodir{../../logo/pdf}
\providecommand{\fmneed}{10}
\def\fm@cursnom{}

% clau del curs, color i color del text damunt del color
\newcommand{\fm@curs}[3]{\@namedef{fm@color@#1}{#2}\@namedef{fm@text@#1}{#3}}
\fm@curs{1eso}{E6B84A}{navy}
\fm@curs{2eso}{3B82F6}{white}
\fm@curs{3eso}{FF7A50}{white}
\fm@curs{4eso-a}{A78BFA}{white}
\fm@curs{4eso-b}{6D28D9}{white}
\fm@curs{1batx-mi}{9B1C1C}{white}
\fm@curs{1batx-ccss}{EF6B6B}{white}
\fm@curs{2batx-mii}{1E6B47}{white}
\fm@curs{2batx-ccss}{2BB3A3}{white}

\newcommand{\fmcurs}[1]{%
  \@ifundefined{fm@color@#1}{\PackageError{fentmates}{Curs desconegut: #1}{}}{}%
  \edef\fm@tmp{\noexpand\definecolor{curs}{HTML}{\@nameuse{fm@color@#1}}%
               \noexpand\colorlet{curstext}{\@nameuse{fm@text@#1}}}\fm@tmp
  \colorlet{cursfosc}{curs!72!black}%
  \def\fm@logo{\fm@logodir/vertical-#1.pdf}%
  \def\fm@marca{\fm@logodir/marca-#1.pdf}}

% ---------- capçalera de la primera pàgina ----------
% #1 matèria (petit), #2 curs, #3 tipus de document i unitat, #4 títol
\newcommand{\capcalera}[4]{%
  \gdef\fm@cursnom{#2}%
  \par\noindent
  \begin{tcolorbox}[colback=curs,colframe=curs,boxrule=0pt,arc=2.2mm,width=3.5cm,height=2.15cm,
      valign=center,left=3mm,right=2mm,top=1mm,bottom=1mm,nobeforeafter,box align=center]
    \color{curstext}{\footnotesize #1}\par\vspace{1pt}{\fontsize{19}{21}\selectfont\bfseries #2\par}%
  \end{tcolorbox}\hspace{5mm}%
  \begin{minipage}[c]{\dimexpr\linewidth-3.5cm-5mm-2.9cm\relax}
    {\color{cursfosc}\small\bfseries\MakeUppercase{#3}}\par\vspace{2pt}
    {\color{navy}\fontsize{20}{23}\selectfont\bfseries #4\par}
  \end{minipage}\hfill
  \begin{minipage}[c]{2.6cm}\raggedleft\includegraphics[width=2.6cm]{\fm@logo}\end{minipage}\par
  \vspace{2pt}}

% ---------- barres de secció del color del curs ----------
\renewcommand{\seccio}[2]{\par\needspace{\fmneed\baselineskip}\bigskip
  \begin{tcolorbox}[colback=curs,colframe=curs,coltext=curstext,arc=2mm,boxrule=0pt,left=3mm,right=3mm,top=1.2mm,bottom=1.2mm]
  {\large\bfseries #1\quad #2}
  \end{tcolorbox}\vspace{-2pt}}

% ---------- estrelles de dificultat del color del curs (a Batxillerat no se'n posen) ----------
\@ifundefined{estrelles}{}{%
  \renewcommand{\estrelles}[1]{{\color{curs}\ifcase#1\or\faStar\faStar[regular]\faStar[regular]\or\faStar\faStar\faStar[regular]\or\faStar\faStar\faStar\fi}}}

% ---------- capçalera i peu de les pàgines ----------
% primera pàgina: logos institucionals al peu
\fancypagestyle{primera}{\fancyhf{}\renewcommand{\headrulewidth}{0pt}
  \fancyfoot[L]{\raisebox{-0.32\height}{\includegraphics[height=0.78cm]{logo-gva.png}}\hspace{5mm}%
    \raisebox{-0.32\height}{\includegraphics[height=0.78cm]{logo-vermellar.png}}}
  \fancyfoot[R]{\scriptsize\color{muted}\href{https://creativecommons.org/licenses/by-nc-sa/4.0/deed.ca}{\ccbyncsa}\ \ CC BY-NC-SA 4.0 · fentmates.org\qquad\small\thepage}}
% la resta: el símbol a l'esquerra i el curs a la dreta, amb el títol al centre com abans
\fancyhead[L]{\raisebox{-0.3\height}{\includegraphics[height=0.62cm]{\fm@marca}}}
\fancyhead[R]{\small\bfseries\color{cursfosc}\fm@cursnom}
\renewcommand{\headrule}{\color{curs!45}\hrule height \headrulewidth}
\makeatother

\fmcurs{4eso-a}

\begin{document}
\thispagestyle{primera}

\capcalera{Matemàtiques A}{4t ESO}{Solucions de la fitxa · Unitat 1}{\raggedright Nombres naturals, enters i fraccionaris}
\vspace{8pt}

\begin{tcolorbox}[colback=soft,colframe=soft,arc=2mm,left=3mm,right=3mm,top=1.5mm,bottom=1.5mm]
\small Intenta resoldre cada exercici abans de mirar-ne la solució. Si no t'ix el mateix resultat, repassa els passos o consulta els apunts interactius de la unitat.
\end{tcolorbox}

% =====================================================
\seccio{1}{Nombres naturals. Operacions combinades}

\ex{1} \begin{tasks}(3)
  \task $2^{5}=\res{32}$
  \task $10^{4}=\res{10000}$
  \task $3^{4}=\res{81}$
  \task $\sqrt{64}=\res{8}$
  \task $\sqrt{121}=\res{11}$
  \task $15^{0}=\res{1}$
\end{tasks}

\ex{1} \begin{tasks}(2)
  \task $347=15\cdot\res{23}+\res{2}$
  \task $1000=24\cdot\res{41}+\res{16}$
\end{tasks}

\ex{1} \begin{tasks}(4)
  \task $18-2\cdot 5+3=\res{11}$
  \task $\left(18-2\right)\cdot 5+3=\res{83}$
  \task $48:8\cdot 2=\res{12}$
  \task $48:\left(8\cdot 2\right)=\res{3}$
\end{tasks}

\ex{2} \begin{tasks}(1)
  \task $3^{2}+4\cdot \sqrt{25}-2^{3}=9+20-8=\res{21}$
  \task $50-\left[18:\left(3^{2}-6\right)\right]\cdot 4=50-\left[18:3\right]\cdot 4=50-6\cdot 4=\res{26}$
  \task $2\cdot \left(5^{2}-3\cdot 4\right)+\sqrt{81}=2\cdot \left(25-12\right)+9=2\cdot 13+9=\res{35}$
  \task $\left(\sqrt{100}+2^{3}\right):6+7=\left(10+8\right):6+7=18:6+7=\res{10}$
\end{tasks}

\ex{2} $100-\left(4\cdot 15+3\cdot 8\right)=100-\left(60+24\right)=100-84=\res{16}$\,€

\ex{3} Ha sumat abans de fer la potència i el producte. $5+3\cdot 2^{2}=5+3\cdot 4=5+12=\res{17}$

\ex{3} Costat: $\sqrt{900}=30$\,m. Tanca: $4\cdot30-4=\res{116\text{ m}}$. Cost: $116\cdot12=\res{1392\text{ €}}$.


% =====================================================
\seccio{2}{Nombres enters. Valor absolut}

\ex{1} \begin{tasks}(5)
  \task $\res{-5}$
  \task $\res{-3}$
  \task $\res{+200}$
  \task $\res{-400}$
  \task $\res{-50}$
\end{tasks}

\ex{1} $\res{-12<-7<-1<0<3<4<9}$

\ex{1} \begin{tasks}(3)
  \task $\lvert -15\rvert =\res{15}$
  \task $\lvert 8\rvert =\res{8}$
  \task $\lvert 0\rvert =\res{0}$
  \task $\text{op}(-6)=\res{6}$
  \task $\text{op}(\lvert -9\rvert )=\res{-9}$
  \task $\lvert -3\rvert +\lvert -5\rvert =\res{8}$
\end{tasks}

\ex{2} \begin{tasks}(1)
  \task $x=\res{-7}$ i $x=\res{7}$
  \task $\res{-3,\ {-2},\ {-1},\ 0,\ 1,\ 2,\ 3}$
  \task $\res{-2,\ {-1},\ 0,\ 1,\ 2}$
\end{tasks}

\ex{2} \begin{tasks}(2)
  \task $5-\left(-8\right)=5+8=\res{13}$\,°C
  \task $5-9=\res{-4}$\,°C
\end{tasks}


% =====================================================
\seccio{3}{Operacions combinades amb enters}

\ex{1} \begin{tasks}(4)
  \task $-8+5=\res{-3}$
  \task $-6-9=\res{-15}$
  \task $7-\left(-4\right)=\res{11}$
  \task $-3+\left(-11\right)=\res{-14}$
  \task $\left(-6\right)\cdot \left(-7\right)=\res{42}$
  \task $8\cdot \left(-9\right)=\res{-72}$
  \task $\left(-56\right):8=\res{-7}$
  \task $\left(-45\right):\left(-9\right)=\res{5}$
\end{tasks}

\ex{1} \begin{tasks}(1)
  \task $-\left(+4\right)+\left(-8\right)-\left(-6\right)+4-1=-4-8+6+4-1=10-13=\res{-3}$
  \task $12-25+3-8-\left(-6\right)=12-25+3-8+6=21-33=\res{-12}$
  \task $-\left(-3\right)+\left(-1\right)-\left(+3\right)-\left(-8\right)=3-1-3+8=11-4=\res{7}$
\end{tasks}

\ex{2} \begin{tasks}(1)
  \task $18:\left(-3\right)+2\cdot \left(-5\right)=-6+\left(-10\right)=\res{-16}$
  \task $\left(-3\right)\cdot \left(2-7\right)+15:\left(-5\right)=\left(-3\right)\cdot \left(-5\right)+\left(-3\right)=15-3=\res{12}$
  \task $4-2\cdot \left[3-\left(6-10\right)\right]=4-2\cdot \left[3+4\right]=4-2\cdot 7=4-14=\res{-10}$
  \task $-20:4\cdot \left(-2\right)=-5\cdot \left(-2\right)=\res{10}$
  \task $25:\left(-5\right)-\left(-3\right)\cdot \left(4-6\right):2=-5-\left(-3\right)\cdot \left(-2\right):2=-5-6:2=-5-3=\res{-8}$
  \task $-6-\left[4-3\cdot \left(2-5\right)\right]=-6-\left[4-3\cdot \left(-3\right)\right]=-6-\left[4+9\right]=-6-13=\res{-19}$
\end{tasks}

\ex{2} \begin{tasks}(1)
  \task $\left(-2\right)^{3}+3^{2}=-8+9=\res{1}$
  \task $-2^{4}+\left(-2\right)^{4}=-16+16=\res{0}$
  \task $5-2\cdot \sqrt{36}+\left(-1\right)^{7}=5-2\cdot 6+\left(-1\right)=5-12-1=\res{-8}$
  \task $\left[\left(-2\right)^{2}-3^{2}\right]^{2}:\left(-5\right)=\left[4-9\right]^{2}:\left(-5\right)=\left(-5\right)^{2}:\left(-5\right)=25:\left(-5\right)=\res{-5}$
  \task $\sqrt{25-9}-\left(-3\right)^{3}:9=\sqrt{16}-\left(-27\right):9=4-\left(-3\right)=4+3=\res{7}$
  \task $\left(-3\right)^{2}\cdot 2-\sqrt{49}:\left(-7\right)=9\cdot 2-7:\left(-7\right)=18-\left(-1\right)=18+1=\res{19}$
\end{tasks}

\ex{2} $350-4\cdot 95+210=350-380+210=\res{180}$\,€. Sí: després dels rebuts el saldo era $350-380=\res{-30}$\,€.

\ex{3} $-150+60-3\cdot 60+45=-150+60-180+45=\res{-225}$\,m

\ex{3} Mitjana: $\left(-3+2-5+0-4+1-5\right):7=\left(-14\right):7=\res{-2}$\,°C. \quad Diferència: $2-\left(-5\right)=\res{7}$\,°C.


% =====================================================
\seccio{4}{Fraccions. Fraccions equivalents}

\ex{1} \begin{tasks}(3)
  \task $\frac{3}{5}\cdot 80=\res{48}$
  \task $\frac{2}{7}\cdot 63=\res{18}$
  \task $\frac{5}{4}\cdot 120=\res{150}$
\end{tasks}

\ex{1} \begin{tasks}[after-item-skip=6pt](2)
  \task $\frac{18}{24}=\frac{18:6}{24:6}=\res{\frac{3}{4}}$
  \task $\frac{40}{64}=\frac{40:8}{64:8}=\res{\frac{5}{8}}$
  \task $\frac{84}{126}=\frac{84:42}{126:42}=\res{\frac{2}{3}}$
  \task $\frac{105}{180}=\frac{105:15}{180:15}=\res{\frac{7}{12}}$
  \task $\frac{90}{126}=\frac{90:18}{126:18}=\res{\frac{5}{7}}$
  \task $\frac{150}{250}=\frac{150:50}{250:50}=\res{\frac{3}{5}}$
\end{tasks}

\ex{1} \begin{tasks}(1)
  \task $\frac{2}{5}=\res{\frac{4}{10}}=\res{\frac{6}{15}}=\res{\frac{8}{20}}$ (per exemple)
  \task $\frac{3\cdot5}{4\cdot5}=\res{\frac{15}{20}}$; \quad $\frac{5\cdot7}{6\cdot7}=\res{\frac{35}{42}}$
  \task $\frac{7\cdot5}{9\cdot5}=\res{\frac{35}{45}}$
\end{tasks}

\ex{2} \begin{tasks}(1)
  \task $6\cdot15=90=9\cdot10$: \res{sí}
  \task $4\cdot20=80\neq7\cdot12=84$: \res{no}
  \task $15\cdot15=225=25\cdot9$: \res{sí}
\end{tasks}

\ex{2} \begin{tasks}(3)
  \task $36x=180\Rightarrow x=\res{5}$
  \task $2x=40\Rightarrow x=\res{20}$
  \task $14x=126\Rightarrow x=\res{9}$
\end{tasks}

\ex{2} Amb denominador comú 12: $\frac{9}{12},\ \frac{8}{12},\ \frac{10}{12},\ \frac{7}{12}$. \quad $\res{\frac{7}{12}<\frac{2}{3}<\frac{3}{4}<\frac{5}{6}}$


% =====================================================
\seccio{5}{Operacions amb fraccions}

\ex{1} \begin{tasks}[after-item-skip=6pt](2)
  \task $\frac{3}{10}+\frac{2}{5}=\frac{3}{10}+\frac{4}{10}=\res{\frac{7}{10}}$
  \task $\frac{5}{6}-\frac{1}{4}=\frac{10}{12}-\frac{3}{12}=\res{\frac{7}{12}}$
  \task $2-\frac{3}{5}=\frac{10}{5}-\frac{3}{5}=\res{\frac{7}{5}}$
  \task $\frac{7}{12}+\frac{5}{18}=\frac{21}{36}+\frac{10}{36}=\res{\frac{31}{36}}$
\end{tasks}

\ex{1} \begin{tasks}[after-item-skip=6pt](2)
  \task $\frac{2}{3}\cdot \frac{5}{7}=\res{\frac{10}{21}}$
  \task $\frac{4}{9}\cdot \frac{3}{8}=\frac{12}{72}=\res{\frac{1}{6}}$
  \task $\frac{3}{5}:\frac{9}{10}=\frac{30}{45}=\res{\frac{2}{3}}$
  \task $6:\frac{3}{4}=\frac{24}{3}=\res{8}$
\end{tasks}

\ex{2} \begin{tasks}[after-item-skip=6pt](2)
  \task $-\frac{3}{4}-\frac{1}{6}+2=-\frac{9}{12}-\frac{2}{12}+\frac{24}{12}=\res{\frac{13}{12}}$
  \task $\left(-\frac{8}{15}\right)\cdot \left(-\frac{25}{12}\right)=\frac{200}{180}=\res{\frac{10}{9}}$
  \task $\frac{7}{10}:\left(-14\right)=-\frac{7}{140}=\res{-\frac{1}{20}}$
  \task $\frac{1}{2}-\frac{2}{3}+\frac{1}{6}=\frac{3}{6}-\frac{4}{6}+\frac{1}{6}=\res{0}$
\end{tasks}

\ex{2} \begin{tasks}(3)
  \task Op. $\res{-\frac{2}{5}}$, \ inv. $\res{\frac{5}{2}}$
  \task Op. $\res{\frac{7}{4}}$, \ inv. $\res{-\frac{4}{7}}$
  \task Op. $\res{-3}$, \ inv. $\res{\frac{1}{3}}$
\end{tasks}


% =====================================================
\seccio{6}{Operacions combinades amb fraccions}

\ex{1} \begin{tasks}[after-item-skip=6pt](2)
  \task $\frac{1}{2}+\frac{1}{3}\cdot \frac{3}{4}=\frac{1}{2}+\frac{1}{4}=\res{\frac{3}{4}}$
  \task $\left(\frac{1}{2}+\frac{1}{3}\right)\cdot \frac{3}{4}=\frac{5}{6}\cdot \frac{3}{4}=\res{\frac{5}{8}}$
  \task $1-\frac{2}{3}:\frac{4}{9}=1-\frac{3}{2}=\res{-\frac{1}{2}}$
  \task $\frac{3}{5}+\frac{2}{5}\cdot \frac{10}{3}=\frac{3}{5}+\frac{4}{3}=\frac{9}{15}+\frac{20}{15}=\res{\frac{29}{15}}$
\end{tasks}

\ex{2} \begin{tasks}[after-item-skip=6pt](1)
  \task $\frac{2}{5}-\left(\frac{1}{2}-\frac{3}{4}\right)-1=\frac{2}{5}-\left(-\frac{1}{4}\right)-1=\frac{8}{20}+\frac{5}{20}-\frac{20}{20}=\res{-\frac{7}{20}}$
  \task $\frac{4}{3}\cdot \left(\frac{2}{5}+\frac{1}{3}\right)=\frac{4}{3}\cdot \frac{11}{15}=\res{\frac{44}{45}}$
  \task $\left(\frac{3}{4}-1\right)\cdot 8+\frac{1}{2}=\left(-\frac{1}{4}\right)\cdot 8+\frac{1}{2}=-2+\frac{1}{2}=\res{-\frac{3}{2}}$
  \task $2-\left[\frac{1}{3}+\frac{1}{4}\right]:\frac{5}{6}=2-\frac{7}{12}:\frac{5}{6}=2-\frac{7}{10}=\res{\frac{13}{10}}$
\end{tasks}

\ex{2} \begin{tasks}[after-item-skip=6pt](1)
  \task $\left(\frac{2}{3}\right)^{2}-\frac{1}{3}\cdot \left(2-\frac{1}{2}\right)=\frac{4}{9}-\frac{1}{3}\cdot \frac{3}{2}=\frac{4}{9}-\frac{1}{2}=\res{-\frac{1}{18}}$
  \task $\frac{5}{6}:\left(\frac{1}{2}-\frac{1}{3}\right)-2^{2}=\frac{5}{6}:\frac{1}{6}-4=5-4=\res{1}$
  \task $\left(-\frac{1}{2}\right)^{3}+\frac{3}{4}:\left(1-\frac{1}{3}\right)=-\frac{1}{8}+\frac{3}{4}:\frac{2}{3}=-\frac{1}{8}+\frac{9}{8}=\res{1}$
\end{tasks}

\ex{3} \begin{tasks}[after-item-skip=6pt](1)
  \task $\left[\frac{3}{4}-\left(\frac{1}{2}\right)^{2}\right]^{2}:\left(-\frac{1}{2}\right)^{3}+\frac{5}{3}\cdot \frac{6}{5}=\left[\frac{3}{4}-\frac{1}{4}\right]^{2}:\left(-\frac{1}{8}\right)+2=\left(\frac{1}{2}\right)^{2}:\left(-\frac{1}{8}\right)+2=\frac{1}{4}:\left(-\frac{1}{8}\right)+2=-2+2=\res{0}$
  \task $\left(\frac{1}{4}+\frac{2}{3}-\frac{1}{12}\right)\cdot 3+\left(\frac{3}{2}-\frac{5}{4}\right):\frac{1}{2}=\frac{10}{12}\cdot 3+\frac{1}{4}:\frac{1}{2}=\frac{5}{2}+\frac{1}{2}=\res{3}$
\end{tasks}


% =====================================================
\seccio{7}{Problemes amb fraccions}

\ex{1} $\frac{2}{5}\cdot30=\res{12}$ amb autobús; \quad $30-12=\res{18}$ no.

\ex{1} $120:\frac{3}{8}=\frac{120\cdot8}{3}=\res{320}$ pàgines.

\ex{2} Lloguer $\res{500}$\,€, menjar $\res{300}$\,€, transport $\res{150}$\,€. \ $1-\left(\frac{1}{3}+\frac{1}{5}+\frac{1}{10}\right)=1-\frac{19}{30}=\res{\frac{11}{30}}$; \ $\frac{11}{30}\cdot1500=\res{550}$\,€.

\ex{2} Dilluns: $\frac{1}{4}\cdot600=150$ L, queden 450 L. \ Dimarts: $\frac{2}{3}\cdot450=300$ L. \ Queden $\res{150\text{ L}}$.

\ex{2} Per persona: $\frac{3}{4}:6=\frac{1}{8}$ kg. \ Per a 10: $10\cdot\frac{1}{8}=\res{\frac{5}{4}\text{ kg}}=1{,}25$ kg.

\ex{3} Després de la roba li queden $\frac{5}{8}$; en el concert gasta $\frac{2}{5}\cdot\frac{5}{8}=\frac{1}{4}$. Ha gastat $\frac{3}{8}+\frac{1}{4}=\frac{5}{8}$ i li queden $\frac{3}{8}$, que són 30\,€: \ $30:\frac{3}{8}=\res{80\text{ €}}$.

\ex{3} Falta $\frac{1}{3}$ de $\frac{3}{4}$, és a dir, $\frac{1}{3}\cdot\frac{3}{4}=\frac{1}{4}$ de l\'etapa: \ $15:\frac{1}{4}=\res{60\text{ km}}$.

\ex{3} Els $\frac{2}{3}$ de l\'oli pesen $1{,}5-0{,}7=0{,}8$ kg; tot l\'oli, $0{,}8:\frac{2}{3}=1{,}2$ kg. Garrafa: $1{,}5-1{,}2=\res{0{,}3\text{ kg}}$.


% =====================================================
\seccio{8}{Potències amb base entera i racional}

\ex{1} \begin{tasks}(3)
  \task $\left(-3\right)^{2}=\res{9}$
  \task $-3^{2}=\res{-9}$
  \task $\left(-2\right)^{5}=\res{-32}$
  \task $\left(-1\right)^{14}=\res{1}$
  \task $\left(-10\right)^{3}=\res{-1000}$
  \task $7^{0}=\res{1}$
\end{tasks}

\ex{1} \begin{tasks}[after-item-skip=6pt](3)
  \task $\left(\frac{2}{3}\right)^{3}=\res{\frac{8}{27}}$
  \task $\left(-\frac{1}{2}\right)^{4}=\res{\frac{1}{16}}$
  \task $\left(-\frac{3}{5}\right)^{3}=\res{-\frac{27}{125}}$
  \task $\left(\frac{5}{4}\right)^{2}=\res{\frac{25}{16}}$
  \task $\left(-\frac{4}{7}\right)^{1}=\res{-\frac{4}{7}}$
  \task $\left(\frac{9}{11}\right)^{0}=\res{1}$
\end{tasks}

\ex{2} \begin{tasks}(4)
  \task $\res{-}$
  \task $\res{+}$
  \task $\res{-}$
  \task $\res{-}$
\end{tasks}

\ex{2} \begin{tasks}[after-item-skip=6pt](2)
  \task $\left(-2\right)^{4}-2^{4}+\left(-2\right)^{3}=16-16+\left(-8\right)=\res{-8}$
  \task $3\cdot \left(-1\right)^{5}-\left(-2\right)^{2}=3\cdot \left(-1\right)-4=-3-4=\res{-7}$
  \task $\left(\frac{1}{2}\right)^{2}+\left(\frac{1}{2}\right)^{3}=\frac{1}{4}+\frac{1}{8}=\frac{2}{8}+\frac{1}{8}=\res{\frac{3}{8}}$
  \task $1-\left(-\frac{2}{3}\right)^{2}=1-\frac{4}{9}=\res{\frac{5}{9}}$
\end{tasks}

\ex{3} $9\cdot \left(\frac{2}{3}\right)^{3}=9\cdot \frac{8}{27}=\res{\frac{8}{3}}$\,m $\approx2{,}67$\,m


% =====================================================
\seccio{9}{Propietats de les potències}

\ex{1} \begin{tasks}(3)
  \task $5^{3}\cdot 5^{4}=\res{5^{7}}$
  \task $2^{9}:2^{4}=\res{2^{5}}$
  \task $\left(3^{2}\right)^{5}=\res{3^{10}}$
  \task $4^{3}\cdot 5^{3}=\res{20^{3}}$
  \task $18^{6}:6^{6}=\res{3^{6}}$
  \task $\left(-7\right)^{4}\cdot \left(-7\right)^{3}=\res{(-7)^{7}}$
\end{tasks}

\ex{1} \begin{tasks}[after-item-skip=6pt](2)
  \task $\left(\frac{2}{5}\right)^3\cdot\left(\frac{2}{5}\right)^4=\res{\left(\frac{2}{5}\right)^{7}}$
  \task $\left(\frac{3}{4}\right)^8:\left(\frac{3}{4}\right)^5=\res{\left(\frac{3}{4}\right)^{3}}$
  \task $\left[\left(-\frac{1}{2}\right)^3\right]^2=\res{\left(-\frac{1}{2}\right)^{6}}$
  \task $\left(\frac{2}{3}\right)^4\cdot\left(\frac{9}{4}\right)^4=\res{\left(\frac{3}{2}\right)^{4}}$
\end{tasks}

\ex{2} \begin{tasks}[after-item-skip=6pt](1)
  \task $\left(-2\right)^{7}:\left(-2\right)^{4}=\left(-2\right)^{3}=\res{-8}$
  \task $\left[\left(-1\right)^{3}\right]^{5}\cdot \left(-1\right)^{2}=\left(-1\right)^{15}\cdot \left(-1\right)^{2}=\left(-1\right)^{17}=\res{-1}$
  \task $2^{5}\cdot 2^{3}:2^{6}=2^{8}:2^{6}=2^{2}=\res{4}$
  \task $\left(\frac{2}{3}\right)^{5}\cdot \left(\frac{2}{3}\right)^{2}:\left(\frac{2}{3}\right)^{4}=\left(\frac{2}{3}\right)^{7}:\left(\frac{2}{3}\right)^{4}=\left(\frac{2}{3}\right)^{3}=\res{\frac{8}{27}}$
  \task $\left(-\frac{1}{3}\right)^{6}:\left(-\frac{1}{3}\right)^{4}=\left(-\frac{1}{3}\right)^{2}=\res{\frac{1}{9}}$
\end{tasks}

\ex{2} \begin{tasks}(1)
  \task \res{F}: $2^3+2^4=8+16=24$
  \task \res{V}
  \task \res{F}: $(2+3)^2=25$, \ $2^2+3^2=13$
  \task \res{V}: exponent parell
  \task \res{F}: $5^0=1$
\end{tasks}

\ex{3} \begin{tasks}[after-item-skip=6pt](1)
  \task $\dfrac{2^4\cdot5^4\cdot2^3}{5^3\cdot2^{10}}=\dfrac{2^7\cdot5^4}{2^{10}\cdot5^3}=\dfrac{5}{2^3}=\res{\frac{5}{8}}$
  \task $3^6\cdot3^3:3^7=3^9:3^7=3^2=\res{9}$
\end{tasks}

\ex{3} $\left(\frac{4}{9}\right)^2=\left[\left(\frac{2}{3}\right)^2\right]^2=\left(\frac{2}{3}\right)^4 \ \Rightarrow\ \left(\frac{2}{3}\cdot\frac{3}{2}\right)^4=1^4=\res{1}$


\end{document}
